\providecommand{\candidate}[1]{}

\section{Introduction}
\label{sec:intro}

On 14 February 2019, a car bomb struck a paramilitary convoy at Pulwama, in
Kashmir, killing more than 40 Indian
personnel~\cite{aljazeera2019pulwama}.
Twelve days later, the Indian Air Force (IAF) struck Balakot, inside
Pakistan~\cite{aljazeera2019balakot}.
On 22 April 2025, three gunmen killed
26 civilians at
Pahalgam~\cite{aljazeera2025pahalgam};
just over two weeks later, IAF launched Operation
Sindoor~\cite{aljazeera2025sindoor}.
Pakistan answered with Operation Bunyan
al-Marsoos~\cite{aljazeera2025bunyan},
and four days of missile, drone and artillery exchanges between two
nuclear-armed states ended in a ceasefire on 10
May~\cite{aljazeera2025ceasefire}. Each of these crises was argued out in public, at the same moment, in
communities belonging to neither government.Reddit hosts a national community for each side. \texttt{r/india} and \texttt{r/pakistan} are the largest general-purpose forums for their respective countries on the platform, each with
its own moderators, its own rules and its own voting public. Alongside them sits \texttt{r/worldnews}, one of the platform's largest international news communities, and one aligned with neither state. A crisis between India and Pakistan arrives in all three at once, and the argument that follows is never one argument.

A cross-border strike is one event, and the communities that discuss it produce more than one narrative~\cite{entman1991kal,tyagi2020ipconflict}. In
one, it arrives as a question of security and necessity; in another, as a question of who was harmed~\cite{makhortykh2017visual}. As the strikes began on 6 May 2025, a comment in \texttt{r/pakistan} drew 153
upvotes for asking readers to ``look after each other and your own families\ldots some of us have never seen war in our
lifetimes.''
% \footnote{\url{reddit.com/comments/1kghif0/\_/mqysggb/}}
Two days later, with the exchanges under way, a comment in \texttt{r/india} drew 182 for putting the question to its own government: ``\ldots Why were the villagers not
evacuated then? Why allow the loss of
life?''
% \footnote{\url{reddit.com/comments/1khoj9j/\_/mr8jwr3/}}
When the ceasefire held on 12 May, a thread in \texttt{r/worldnews} gave 1,151 upvotes to ``[n]o more civilian casualties. There can be no greater
victory than this.''
% \footnote{\url{reddit.com/comments/1kknm0z/\_/mrvszsy/}}
All three carry the same narrative frame, because each makes civilian harm the thing at stake. What each does with it is different, and so is the community it was
posted in. We therefore treat community alignment as a property of the subreddit and not of any speaker's nationality, and the second comment shows why:
an India-aligned forum turning on its own government. Two features separate this setting from the intra-societal hostility most harm research is built
on~\cite{zampieri2019olid,davidson2017hatespeech,salehabadi2022user,aleksandric2025power}. Alignment tracks a territorial dispute instead of elective membership~\cite{fang2022indivisibility}, and the clock that drives the discourse runs outside the platform.


% OLD:
% Conflict discourse has been studied closely in CSCW and HCI, and mostly one event at a time. Crisis informatics established what platform traces reveal while a
% disaster unfolds~\cite{vieweg2010microblogging} and built a discipline for comparing communications across crises~\cite{olteanu2015crisis}, while remaining
% reflexive about how its own instruments bound what it can see~\cite{soden2018informating}. \citet{stewart2017lines} traced networked frame
% contests inside a single protest movement, where rival framings compete within one discourse instead of across two, and \citet{zhou2025harm} composed misinformation,
% narrative and target annotations over anti-Asian speech, then measured how those content factors and the judge's own background bear on how harmful the speech is
% perceived to be. The construct those studies turn on was given its standard definition by \citet{entman1993framing}, who describes framing as selecting ``some aspects of a perceived reality'' and making them ``more salient in a communicating text, in
% such a way as to promote a particular problem definition, causal interpretation, moral evaluation, and/or treatment recommendation.''
% A narrative frame here is that selection. It is the explanation a comment advances for the conflict, not the subject it mentions: a comment about an airstrike that turns on who was
% killed is humanitarian, while the same strike discussed as capability and deterrence is security and military. We use six such frames. Computational work
% has carried frame inventories onto social
% platforms~\cite{card2015mediaframes,mendelsohn2021framing} and modelled how frames and events move together over time~\cite{hopp2020markov}. On interstate
% dyads, \citet{shin2026diplomacy} separate framing from affect around a diplomatic shock, finding that Reddit sentiment toward North Korea reverted once the Hanoi
% summit failed, while the shift away from threat-oriented framing did not.
Conflict discourse has been studied closely in CSCW and HCI, and mostly one
event at a time. \citet{stewart2017lines} traced networked frame contests inside
a single protest movement, where rival framings compete within one discourse
instead of across two, and \citet{zhou2025harm} composed misinformation,
narrative and target annotations over anti-Asian speech. Where this work has
compared across events rather than within one, it has compared separate
crises rather than the communities arguing about a single
one~\cite{olteanu2015crisis}. The construct these studies turn on is the frame,
which \citet{entman1993framing} defines as selecting ``some aspects of a
perceived reality'' and making them ``more salient'', so as to promote a
particular problem definition, causal interpretation, moral evaluation or
treatment recommendation. A narrative frame here is that selection: the
explanation a comment advances for the conflict, not the subject it mentions. A
comment about an airstrike that turns on who was killed is humanitarian, while
the same strike discussed as capability and deterrence is security and military.
We use six such frames. Computational work has carried frame inventories onto
social platforms and modelled how frames and events move together over
time~\cite{card2015mediaframes,mendelsohn2021framing}. On interstate dyads,
\citet{shin2026diplomacy} separate framing from affect around a diplomatic
shock, finding that Reddit sentiment toward North Korea reverted once the Hanoi
summit failed, while the shift away from threat-oriented framing did not.


The three comments above are one frame in three audiences. Whether that difference is systematic is not a question this literature asks, and the reason
is structural. Harm and stance are labelled one construct at a time because separate benchmarks are what annotation budgets and pretrained classifiers
afford~\cite{zampieri2019olid,sap2020socialbias}, and the compositional work that exists composes inside the content layer~\cite{zhou2025harm,scheuerman2021severity}. Where context has been admitted, it has been the wrong context. \citet{pavlopoulos2020context} ask whether the surrounding conversation bears on a comment's toxicity label, leaving the community itself fixed. Yet norms on Reddit are stratified rather than
platform-wide~\cite{chandrasekharan2018hiddenrules}, and communities differ
systematically in the affect they express and
reward~\cite{waller2021quantifying,rathje2021outgroup}. The same gap runs through work on boundary-crossing, reported as friction in some
accounts~\cite{kumar2018conflict,garimella2018echo,bail2018exposure} and as workable under the right design in
others~\cite{rajadesingan2021crosspartisan,xia2025integrated}. \mohit{Here say that the prior works on the Indo-Pak conflict}
Computational work on India--Pakistan conflict discourse has concentrated on the
weeks around one event, the February 2019 Pulwama attack and the Balakot strike
that answered it. \citet{tyagi2020ipconflict} infer stance from hashtags on
Indian and Pakistani Twitter across both. \citet{hussain2021pulwama} hand-codes
tweets from the attack into four categories and splits them by national side.
\citet{palakodety2020hope} classify hope speech in code-mixed YouTube comments
over the same pair of events. Despite the scale at which this conflict is argued
online, little research has gone to how one crisis is rendered differently by the
communities receiving it: each of these studies covers a single crisis with one
or two label layers, and none of them models the community a comment appeared
in. The nearest Reddit study of an adjacent conflict annotates a single stance
layer over a corpus roughly thirty times smaller than
ours~\cite{ali2025stance}.
\yash{Named the three prior Indo-Pak studies (Tyagi, Hussain, Palakodety) with
platform, event and label depth. Rewrote the lead-in to name the dyad and the
event explicitly instead of ``this dyad'' and ``it'', and moved the delta into an
explicit gap sentence.}



\mohit{We should mention in small how we collected the data, say using Arcticshift API, we analyse ... }
\mohit{Provide the names of them in bracket so after you finish six India alligned put the names}
% OLD:
% To examine how narrative frame, community, and crisis type combine, we analyse 305,217 comments posted to ten subreddits across six India--Pakistan crises between January 2019 and January 2026: six India-aligned communities, a Pakistan-aligned community, and the three general-audience communities in which this conflict is argued at scale.
To examine how narrative frame, community and crisis type combine, we retrieved
Reddit discussion of six India--Pakistan crises from the Arctic Shift
archive~\cite{arcticshift2026software} through the BAScraper
client~\cite{bascraper2026software}, and we analyse the 305,217 comments posted
to ten communities between January 2019 and June 2025. Six of those communities
are India-aligned (\texttt{r/india}, \texttt{r/indianews},
\texttt{r/IndiaSpeaks}, \texttt{r/bakchodi}, \texttt{r/IndianDefense} and
\texttt{r/indiadiscussion}), one is Pakistan-aligned (\texttt{r/pakistan}), and
three are the general-audience communities in which this conflict is argued at
scale (\texttt{r/worldnews}, \texttt{r/geopolitics} and
\texttt{r/lesscredibledefence}).
\yash{Added the Arctic Shift + BAScraper retrieval route in one clause, and
named all ten communities.}
\mohit{Say which one and start with to investigate this multiyear something like that we used a...}
% OLD:
% Four of the six events are kinetic military events forming two attack-and-retaliation pairs. The other two are the 2019 constitutional change in Kashmir and a 2021 ceasefire agreement, so the design separates the kind of crisis from the year it happened.
To investigate a conflict that keeps returning, we built a panel that separates
the kind of crisis from the year it occurred. Four of the six events are
military, forming two attack-and-retaliation pairs six years apart: the Pulwama
attack and the Balakot strike in February 2019, and the Pahalgam attack and
Operation Sindoor in April and May 2025. The remaining two events are
non-military, the revocation of Article 370 in August 2019 and the ceasefire
agreement of February 2021.
\yash{Opened on the multi-year design as asked, and named which events are
military and which are not, with dates.}
% OLD:
% A large language model labels each comment for sentiment, stance toward each
% state and narrative frame, and a toxicity classifier supplies the offensiveness label, so a comment can be read as a composition instead of a single score.
Each comment carries seven annotation layers. A commercial large language model
labelled content type, narrative frame, sentiment and stance toward each state,
and a toxicity classifier supplied the offensiveness label, so a comment can be
read as a composition instead of a single score. We chose that labeller by
measurement rather than by reputation, scoring four families of system against
human gold standards before adopting any of them, and every machine agreement
figure in \S\ref{sec:data} appears beside the human figure for the same
construct. \mohit{You should give more details about the methods}
\yash{Added the layer count, the labeller-selection design and the paired
human/machine reporting standard, keeping the protocol itself in \S3.}
% Section~3 reports agreement against double-coded human gold standards for every layer.
We measure contestation with Reddit's \texttt{controversiality} flag, which marks
comments drawing substantial votes both ways. \mohit{You need some kind of bridge sentence here, not the one that you had. Also, there needs to be a sentence maybe earlier somewhere where we say, to the best of our knowledge this is the first longitudinal study .....}
To the best of our knowledge, this is the first longitudinal study of
India--Pakistan crisis discourse to span more than one crisis, and the first to
model the narrative frame a comment advances together with the community that
received it. We therefore treat a comment as a pairing of the frame it carries
with the community it was posted to, and ask what that pairing explains.
\yash{Added the primacy claim, scoped to two narrow versions that the design
supports, and replaced the jump into the RQs with a bridge sentence that states
the frame-community pairing the questions then test.}
% The analysis is primarily quantitative, with qualitative close
% reading used to interpret the patterns it surfaces.
In particular, we investigate the following research questions:

\begin{itemize}
\item \textbf{RQ1:} Which combinations of narrative frame and community make crisis discourse contested, and how do those combinations compare with the content-layer toxicity signal that harm research typically measures?
\item \textbf{RQ2:} Why do those combinations differ, and is an event's frame--affect signature fixed by the kind of crisis or by the year it occurred?
\item \textbf{RQ3:} What happens at the community boundary, where participants cross between opposed communities?
\end{itemize}



Our results show that what a community argues over turns on the frame a comment advances and the community it appears in, and hardly at all on how offensively it is worded. Humanitarian framing is the most contested category at 10.4\% and economic framing the least at 5.9\%; across the three alignment groups the range is wider, from 11.7\% in the general-audience communities to 4.8\% in the India-aligned ones. Community alignment separates controversiality rates 2.5-fold and narrative frame 1.8-fold, against 1.12-fold for the offensive-versus-not contrast that content moderation acts on. Measured against a toxicity threshold, offensive and inoffensive comments differ by 0.8 points, and inside humanitarian framing they barely differ at all, 10.2\% against 10.5\%; in a subreddit fixed-effects logit the frame gradient survives absorbing community composition while the offensive coefficient stays negligible (odds ratio $1.17$, $p = 3.4\times10^{-11}$). What divides a community is the claim, not the tone in which it arrives.
% OLD:
% The same frame is also not the same object in two communities. Security and military framing is rendered positively
% about twice as often in India-aligned communities as in general-audience ones, an interaction that holds once stance toward both states is controlled, and the
% communities reward it at different rates. An event's signature, meaning the mix of frames used and the sentiment attached to them, tracks the kind of crisis more closely than the year: \textit{Pulwama} and \textit{Operation Sindoor} lie six years apart and are
% close in composition, while Pulwama and the constitutional change six months later are far apart. That separation rests on interval estimates. At the
% community boundary, crossing intensifies opposition instead of softening it. No discrete bridge-user type appears in these data, those who cross are drawn
% disproportionately from the more hawkish participants, and tracking 1,606 authors against themselves shows the same people expressing more anti-out-group stance
% when they post across the line.
The same frame is also not the same object in two communities. Under security and
military framing, the adjusted probability that a comment reads as positive is
0.094 in India-aligned communities against 0.050 in general-audience ones, and
the frame-by-community interaction \emph{strengthens} when stance toward both
states is controlled ($\chi^2(20) = 287.6$, $p = 2.7\times10^{-49}$), so the
asymmetry is not partisanship in another guise. The communities also reward that
affect differently: the same negatively-read military comment earns an adjusted
1.239 in the Pakistan-aligned community against 1.418 in India-aligned
communities and 1.566 in general-audience ones, which makes the asymmetry one
that is produced and then paid for. An event's signature, meaning the mix of
frames used and the sentiment attached to them, tracks the kind of crisis more
closely than the year. \textit{Pulwama} and \textit{Operation Sindoor} lie six
years apart and sit at a cosine distance of 0.0074 (95\% CI [0.0022, 0.0313]),
while Pulwama and the constitutional change six months later sit at 0.2058. At
the community boundary, crossing intensifies opposition instead of softening it.
No discrete bridge-user type appears in these data, those who cross concentrate
among the more hawkish participants ($\chi^2(2) = 301.9$, Cram\'er's $V = 0.226$,
$p = 2.7\times10^{-66}$), and tracking 1,606 authors against themselves shows the same people
raising their anti-out-group stance from 0.339 at home to 0.376 across the line
($+0.037$, $p < .001$, paired Wilcoxon).
\mohit{Provide more stastical answers here, use p values etc. This is a good start. Also dont use \sout{} no ned for it, just comment the old things and put new stuff and then your name comment}
\yash{Added the statistics Jun flagged as mine: the R5 stance-controlled
interaction with its $p$, the adjusted positivity probabilities, the R5 reward
figures, the R2 cosine distances with the interval, the R6 hawkish
concentration, and the R6 paired-Wilcoxon result with $n$ and $p$. All checked
against \S\ref{sec:results}.}
\jun{Added statistical detail to the R8 sentences (my part): the 2.5$\times$/1.8$\times$/1.12$\times$ separations and the subreddit-FE offensive OR $1.17$, $p=3.4\times10^{-11}$; also corrected the within-Humanitarian split to match \S5.6. The R5 interaction ($\chi^2(20)=287.6$, $p=2.7\times10^{-49}$) and the R6 paired-Wilcoxon result ($+0.037$, $p<0.001$) are the natural p-values to add to the affective-privileging and boundary-crossing sentences --- leaving those to Yash since they own R5/R6.}



In each of these results, the community is doing work that a comment-level measure cannot see.
We make three contributions to CSCW, HCI and social computing research on conflict discourse and online communities.
(1) \textit{Contestation tracks frame and community rather than tone} (RQ1). We show that what a
community contests is a property of the frame and the community together, and hardly at all of a comment's tone. Humanitarian framing splits a community even when politely worded. Security and military framing draws it together. How offensive a comment is barely predicts whether a community will fight about it.
(2) \textit{A frame is not one object across communities} (RQ2). We find that the same frame is produced and rewarded differently depending on where it appears, and that an event's
signature tracks the kind of crisis rather than the year. A frame label read without the community it appeared in loses what separates one community from
another.
(3) \textit{Crossing a community boundary is incursion rather than contact} (RQ3). No bridge-user type appears in the data. Those who cross skew
hawkish, and the receiving community penalises civil and hostile outsiders alike. Because we follow the same authors across the boundary, this is a change in how
people write and not in who shows up. These findings bear on how harm is measured in conflict discourse, and on any design that routes people across community lines.
\mohit{Here you need a paragraph that closes it. Say in small what we did and in like 3 lines gives the results and then close with how this work is important in one line }



Across six crises between 2019 and 2025, we labelled 305,217 Reddit comments on seven layers and modelled what each of ten communities fought over. What a community contests turns on the story a comment tells and the community it appears in, and hardly at all on how offensively it is worded. The same frame is produced and rewarded differently on either side of an alignment boundary, and an event's signature follows the kind of crisis rather than the year it
occurred. Crossing that boundary sharpens opposition rather than softening it,
and the receiving community penalises civil and hostile outsiders alike. For
CSCW, our results carry a practical consequence: conflict discourse has to be
measured at the pairing of a comment with the community that received it,
because neither term predicts contestation on its own.
\yash{Added the closing paragraph: one sentence on what we did, three on the
results, one on why it matters to CSCW.}




% \mohit{I am unsure why this is coming here? This should come in the para where we have how much data we collected in the intro, but in short way. This can go in the methods when you are explaining things}We had intended to establish these patterns causally and could not. The crises recur inside each other's recovery windows, so the calm period before one event
% is the aftermath of the one before it, and parallel trends fails on both military events whose pre-period is long enough to test.
% Inference is a second problem. We cluster standard errors by community, which leaves seven to nine clusters
% depending on the event, and at that few the usual cluster-robust estimator
% over-rejects. It marked nine of eighteen per-event estimates as significant and a
% wild cluster bootstrap left
% none~\cite{cameron2008bootstrap,mackinnon2017wildbootstrap}. The design was also
% underpowered: no cell reached its minimum detectable effect, so the bootstrap
% null means this panel cannot detect an effect, not that none exists. On a panel of this
% shape the design does not identify the effects asked of it, so we report a
% compositional account instead.